\documentclass[10pt,a4paper]{article}
\usepackage{fontspec,amsmath,amssymb,graphicx}
\usepackage[margin=18mm]{geometry}
\setmainfont{Times New Roman}
\setlength{\parindent}{0pt}
\setlength{\parskip}{5pt}
\allowdisplaybreaks
\sloppy
\emergencystretch=2em
\begin{document}
\section*{BT\_5\_6\_3\_04 — Lời giải bổ sung}
(a) Đặt $\phi=e^{y-x}v$. Khi đó $(\Delta+2\partial_x-2\partial_y)\phi=e^{y-x}(\Delta v-2v)$. Do đó $\phi_{nm}=e^{y-x}\sin n\pi x\sin m\pi y$, $\lambda_{nm}=\pi^2(n^2+m^2)+2$, $n,m\ge1$.
(b) Chọn $\sigma=e^{2x-2y}$, tích có trọng số biến thành tích các sine, nên trực giao.
(c) $u=3+e^{y-x}\sum_{n,m\ge1}c_{nm}e^{-\lambda_{nm}t}\sin n\pi x\sin m\pi y$, với $c_{nm}=4(-1)^{n+m}/(nm\pi^2)$ là hệ số sine của $xy$.
(d) Mọi $\lambda_{nm}>0$, nên giới hạn bằng 3. Dữ liệu đầu không khớp biên ở cạnh phải/trên; nghiệm hiểu theo dữ liệu nội và biên cho $t>0$.\end{document}